\chapter{Your Thesis as a Hardware Workload}\label{ch:workload}

Before you design any hardware, you have to know exactly what you are asking it
to do: which operations, how many per second, how much memory must persist
between frames, and under what deadline. Hardware designers call this
\emph{characterising the workload}. It is the bridge between your thesis and
everything that follows. Without it, design decisions are guesses. With it,
most decisions in Phases~1 and~2 follow almost mechanically.

This chapter characterises the thesis system~\cite{burrell2026thesis}. Every
number marked \emph{measured} comes from
\file{docs/research_guide/scripts/measure_workload.py}, which instruments the
thesis code directly. You will rerun it yourself in \cref{st:workload:reproduce}.
Numbers marked \emph{estimate} are derived by hand from the code and are only as
good as the counting. \Cref{app:factsheet} collects all of them on one page.

\begin{whybox}
Most student hardware projects fail quietly here. They start writing RTL before
they know where the cost is, then optimise the wrong block. Your thesis has an
unusual advantage: its code is short, readable and already organised into the
same blocks a chip would have. This chapter turns that code into a cost model.
\end{whybox}

% =========================================================================
\section{Two planes, two rates}\label{sec:workload:planes}

Your thesis already has the split that hardware designers use to organise real
systems (\cref{fig:workload:planes}). The \emph{data plane} is the audio path.
It processes every sample, never stops, and has a hard deadline: if a frame is
late, the user hears a glitch. The \emph{control plane} is the intent path. It
runs only when the user says something, and its only output is a handful of
parameters.

\begin{figure}[htbp]
\centering
\begin{tikzpicture}[
  font=\sffamily\footnotesize,
  blk/.style={draw=#1,fill=#1!8,rounded corners=1.2mm,minimum height=8mm,
              minimum width=17mm,align=center,line width=0.6pt},
  blk/.default=brand,
  arr/.style={-{Stealth[length=2mm]},line width=0.7pt},
  ctl/.style={-{Stealth[length=2mm]},line width=0.7pt,draw=pitc,dashed},
  node distance=4mm]
  % data plane
  \node[blk] (mic) {Mic\\\num{48}\,kHz};
  \node[blk,right=of mic] (stft) {STFT\\(window+FFT)};
  \node[blk,right=of stft] (imcra) {IMCRA\\noise PSD};
  \node[blk,right=of imcra] (wien) {Parametric\\Wiener};
  \node[blk,right=of wien] (istft) {ISTFT\\(IFFT+OLA)};
  \node[blk,right=of istft] (spk) {Receiver};
  \draw[arr] (mic) -- (stft); \draw[arr] (stft) -- (imcra);
  \draw[arr] (imcra) -- (wien); \draw[arr] (wien) -- (istft); \draw[arr] (istft) -- (spk);
  \begin{scope}[on background layer]
    \node[draw=brand,fill=brandlight!60,rounded corners=2mm,inner sep=3mm,
          fit=(mic)(spk),label={[anchor=south west,font=\sffamily\bfseries\footnotesize,text=brand]north west:Data plane --- every frame, hard deadline \qty{10.667}{\ms}}] {};
  \end{scope}
  % control plane
  \node[blk=rqc,below=17mm of mic] (txt) {Complaint\\(text)};
  \node[blk=rqc,right=of txt] (emb) {Sentence\\encoder};
  \node[blk=rqc,right=of emb] (cls) {Intent class\\+ severity};
  \node[blk=rqc,right=of cls] (map) {Parameter\\mapping};
  \draw[arr] (txt) -- (emb); \draw[arr] (emb) -- (cls); \draw[arr] (cls) -- (map);
  \begin{scope}[on background layer]
    \node[draw=rqc,fill=rqc!4,rounded corners=2mm,inner sep=3mm,
          fit=(txt)(map),label={[anchor=north west,font=\sffamily\bfseries\footnotesize,text=rqc]south west:Control plane --- only when the user speaks}] {};
  \end{scope}
  \draw[ctl] (map.north) -- node[left,align=right,text=pitc,font=\sffamily\scriptsize]{$\beta$, floor\\$\Rightarrow$ gain LUT} (map.north |- wien.south);
  \draw[ctl] (map.east) -| node[pos=0.75,right,align=left,text=pitc,font=\sffamily\scriptsize]{output gain, tilt\\$\Rightarrow$ 2 registers} (istft.south);
\end{tikzpicture}
\caption{The thesis system redrawn as a data plane and a control plane. The
dashed red arrows are the \emph{entire} interface between them. In hardware it
shrinks to a small lookup table and two registers (opportunities O1 and O2 in
\cref{sec:workload:opportunities}).}
\label{fig:workload:planes}
\end{figure}

\begin{table}[htbp]
\centering\small
\caption{The two planes have completely different hardware requirements.}
\label{tab:workload:planes}
\begin{tabularx}{\textwidth}{@{}lLL@{}}
\toprule
 & \textbf{Data plane (audio)} & \textbf{Control plane (intent)} \\
\midrule
Rate            & \qty{93.75}{\Hz} frames, \qty{48}{\kHz} samples & once per utterance (seconds to minutes apart) \\
Deadline        & hard: \qty{10.667}{\ms} per frame, every frame & soft: a few hundred milliseconds feels instant \\
Dominant cost   & memory for per-bin state; divisions and transcendental functions & the sentence encoder (\num{e8}--\num{e9} operations per utterance) \\
Right metric    & energy per frame, area, latency & energy per inference, model size \\
Thesis code     & \file{src/intent_se/audio/} & \file{src/intent_se/nlp/}, \file{src/intent_se/control/} \\
Hardware style  & fixed-function streaming datapath & small accelerator, or software on a CPU or phone \\
\bottomrule
\end{tabularx}
\end{table}

% =========================================================================
\section{The frame budget}\label{sec:workload:budget}

The hop length $H$ fixes how often a frame must be produced. With
$f_s=\qty{48}{\kHz}$ and $H=512$ (\file{src/intent_se/config.py}):
\begin{equation}
  f_\text{frame} = \frac{f_s}{H} = \qty{93.75}{\Hz},
  \qquad
  T_\text{frame} = \frac{1}{f_\text{frame}} = \qty{10.667}{\ms}.
  \label{eq:workload:frame}
\end{equation}
A chip running at clock frequency $f_\text{clk}$ therefore has
$N_\text{cyc}=f_\text{clk}\,T_\text{frame}$ clock cycles to finish each frame
(\cref{tab:workload:cycles}). That cycle count is the most important number in
Phase~2. It tells you how much hardware you can \emph{share}: if a frame needs
\num{100000} operations and you have \num{1000000} cycles, a single arithmetic
unit can do everything, ten times over.

\begin{table}[htbp]
\centering\small
\caption{Clock cycles available per frame (measured frame period).}
\label{tab:workload:cycles}
\begin{tabular}{@{}rr@{}}
\toprule
$f_\text{clk}$ & cycles per frame \\
\midrule
\qty{100}{\MHz} & \num{1066667} \\
\qty{50}{\MHz}  & \num{533333} \\
\qty{10}{\MHz}  & \num{106667} \\
\qty{1}{\MHz}   & \num{10667} \\
\bottomrule
\end{tabular}
\end{table}

\begin{conceptbox}{Throughput is easy; energy is the problem}
The thesis code already runs in about \qty{0.2}{\ms} per frame in Python on an
ordinary computer (measured: real-time factor $\approx 0.02$). Speed is not the
challenge. The challenge is doing the same work in a hearing-aid power budget.
Dynamic power scales as $\alpha C V^2 f$~\cite{chandrakasan1992lowpower}, so the
winning strategy is usually to run \emph{slowly} and at low voltage, and to touch
memory as little as possible. Keep that in mind whenever this guide asks
\enquote{how few cycles can we use?}: the real question is \enquote{how little
energy per frame?}
\end{conceptbox}

% =========================================================================
\section{Operation inventory of the data plane}\label{sec:workload:ops}

\Cref{tab:workload:ops} walks through one frame of the data plane, stage by
stage, in the order the code runs it. The per-bin counts come from reading the
code as written, with no algebraic simplification yet. Treat them as estimates
accurate to within a few operations. Two things stand out. First, \emph{every}
stage after the FFT is per-bin and independent across bins, apart from the
5-tap frequency smoothing. That makes it ideal for a streaming pipeline that
processes one bin per clock cycle. Second, the expensive operations are
divisions (12 per bin) and two transcendental functions, not
multiplications.

\begin{table}[htbp]
\centering\footnotesize
\setlength{\tabcolsep}{4pt}
\caption{Per-bin operations per frame, counted from the code (estimate).
Columns: multiplications, additions/subtractions, divisions, comparisons
(including \code{min}, \code{max}, \code{clip}), and transcendental functions.
Line numbers refer to the thesis repository.}
\label{tab:workload:ops}
\begin{tabularx}{\textwidth}{@{}Ll*{5}{c}@{}}
\toprule
Stage & Source & $\times$ & $\pm$ & $\div$ & cmp & exp/pow \\
\midrule
Power spectrum $|X_k|^2$ & \file{pipeline.py:179} & 2 & 1 & -- & -- & -- \\
Frequency smoothing (5 taps) & \file{imcra.py:117} & 5 & 4 & -- & -- & -- \\
Recursive smoothing $S$ & \file{imcra.py:118} & 2 & 1 & -- & -- & -- \\
Minimum tracker 1 & \file{imcra.py:40-51} & -- & -- & -- & 9 & -- \\
$\gamma_\text{min}$, $\zeta$ (shared denominator) & \file{imcra.py:121-123} & 1 & -- & 2 & 1 & -- \\
Noise-only indicator & \file{imcra.py:126} & -- & -- & -- & 2 & -- \\
Second smoothing $\tilde S_f$ & \file{imcra.py:129-131} & 6 & 8 & 1 & 1 & -- \\
Recursive smoothing $\tilde S$ & \file{imcra.py:133} & 2 & 1 & -- & -- & -- \\
Minimum tracker 2 & \file{imcra.py:134} & -- & -- & -- & 9 & -- \\
$\tilde\gamma_\text{min}$, $\tilde\zeta$ (shared denominator) & \file{imcra.py:136-138} & 1 & -- & 2 & 1 & -- \\
Speech-absence probability $q$ & \file{imcra.py:141-146} & 1 & 1 & -- & 6 & -- \\
A posteriori SNR $\gamma$, $\nu$ & \file{imcra.py:149-152} & 1 & 1 & 2 & 1 & -- \\
Likelihood ratio, presence prob.\ $p$ & \file{imcra.py:153-154} & 2 & 3 & 2 & 2 & exp \\
Noise PSD update $\lambda_d$ & \file{imcra.py:159-164} & 4 & 3 & -- & 1 & -- \\
Wiener: $\gamma$, decision-directed $\xi$ & \file{wiener.py:89-98} & 2 & 2 & 2 & 3 & -- \\
Wiener: gain, floor, tilt, clip & \file{wiener.py:102-112} & 2 & 1 & 1 & 3 & pow (+$10^x$) \\
DD memory, apply gain & \file{wiener.py:115}, \file{pipeline.py:196} & 4 & -- & -- & -- & -- \\
\midrule
\textbf{Total per bin} & & \textbf{35} & \textbf{26} & \textbf{12} & \textbf{39} & \textbf{2 (+1)} \\
\bottomrule
\end{tabularx}
\end{table}

Adding the transforms gives a per-frame total (\cref{tab:workload:frameops}).
The real 2048-point FFT can be computed as a 1024-point complex FFT plus a
post-processing pass. A radix-2 complex FFT of size $N$ needs
$\frac{N}{2}\log_2 N = \num{5120}$ butterflies for $N=1024$, and each butterfly
is roughly 10 real operations~\cite{cooley1965fft,oppenheim2010dtsp}.

\begin{table}[htbp]
\centering\small
\caption{Operations per frame and per second (estimate).}
\label{tab:workload:frameops}
\begin{tabular}{@{}lrr@{}}
\toprule
Block & ops per frame & share \\
\midrule
Forward real FFT (2048 points) & $\approx\num{56000}$ & 25\,\% \\
Inverse real FFT (2048 points) & $\approx\num{56000}$ & 25\,\% \\
Analysis window, synthesis window, overlap-add & $\approx\num{3000}$ & 1\,\% \\
Per-bin processing ($1025 \times \approx 112$) & $\approx\num{115000}$ & 50\,\% \\
\midrule
\textbf{Total} & $\boldsymbol{\approx\num{230000}}$ & \\
\textbf{Per second} (${}\times\qty{93.75}{\Hz}$) & $\boldsymbol{\approx\num{21e6}}$ & \\
\bottomrule
\end{tabular}
\end{table}

\begin{whybox}[What the inventory tells you]
At \qty{100}{\MHz} you have \num{1066667} cycles for roughly \num{230000}
operations, so throughput is not the constraint. Two architectures bracket the
design space you will explore in \cref{ch:rtl} (both estimates):
\begin{itemize}
  \item \textbf{Fast and lean.} One radix-2 butterfly per cycle plus a fully
  pipelined per-bin unit (one bin per cycle, initiation interval 1) needs about
  $2\times(\num{5120}+512)+1025+\text{overhead}\approx\num{12500}$ cycles per
  frame. That meets the deadline at only $\approx\qty{1.2}{\MHz}$.
  \item \textbf{Fully folded.} Keep the single-butterfly FFT, but process the
  bins with one shared arithmetic unit doing about one operation per cycle
  (about 100 cycles per bin). That needs $\approx\num{115000}$ cycles per frame,
  so $\approx\qty{11}{\MHz}$.
\end{itemize}
Hypothesis H4 (\cref{ch:overview}) says that the shared-resource design meets
the deadline at \qty{5}{\MHz} or less. The interesting research question is which
point between these extremes gives the lowest \emph{energy per frame}.
\end{whybox}

% =========================================================================
\section{Nonlinear functions: where the real hardware cost is}\label{sec:workload:nonlinear}

In Python, a division or a \code{**} costs about the same as an addition. In
hardware they are in different leagues. A multiplier maps to one FPGA DSP slice.
A divider, an \code{exp} or a \code{pow} is either a long iterative circuit or a
lookup table plus arithmetic (\cref{ch:primer}). \Cref{tab:workload:nonlinear}
lists every nonlinear operation in the data plane and what you will replace it
with.

\begin{table}[htbp]
\centering\footnotesize
\caption{Nonlinear operations per bin per frame and their hardware replacements.}
\label{tab:workload:nonlinear}
\begin{tabularx}{\textwidth}{@{}lcLL@{}}
\toprule
Operation & Count & Where & Hardware replacement \\
\midrule
Division & 12 & 9 in IMCRA, 3 in Wiener & Only 9 distinct denominators, so share reciprocals (O4). Caching $H=\xi/(1+\xi)$ (O3) removes one more. Log domain turns most of them into subtractions (O5). \\
$(\xi/(1+\xi))^{\beta}$ & 1 & \file{wiener.py:102} & Gain lookup table indexed by $\xi$ in dB, rebuilt by the control plane when $\beta$ changes (O1). The datapath never computes \code{pow}. \\
$e^{-\nu}$ & 1 & \file{imcra.py:153} & Small LUT or piecewise-linear approximation over $\nu\in[0,16]$ (O7), or remove it entirely with the log-domain rewrite below. \\
$10^{\text{tilt}\cdot f/20}$ & 1 & \file{wiener.py:110} & A geometric sequence across bins: one multiply per bin by a ratio $r$ the control plane writes (O2). \\
$\min$ over 9 values & 2 & minimum trackers & Comparators. Values can be stored in log domain with few bits (O5). \\
\bottomrule
\end{tabularx}
\end{table}

\begin{conceptbox}{An example of an algebraic rewrite for hardware}
The speech-presence probability (\file{imcra.py:152-154}) is
\begin{equation}
  p = \frac{1}{1+\dfrac{q}{1-q}\,(1+\xi)\,e^{-\nu}},
  \qquad \nu = \gamma\,\frac{\xi}{1+\xi}.
\end{equation}
As written, that is three divisions and an exponential. Take logarithms of the
ratio in the denominator and it becomes exactly
\begin{equation}
  p = \sigma\!\Big(\nu - \ln(1+\xi) - \operatorname{logit}(q)\Big),
  \qquad \sigma(z)=\frac{1}{1+e^{-z}},\quad
  \operatorname{logit}(q) = \ln\frac{q}{1-q}.
  \label{eq:workload:sigmoid}
\end{equation}
In hardware, $\ln(1+\xi)$ is a small table indexed by $\xi$ in dB. Conveniently,
that is the same index as the gain LUT. $\operatorname{logit}(q)$ is a small
table over $q$, and $\sigma$ is one more small table. Three divisions and an
\code{exp} become three table lookups and two additions. The rewrite is exact in
real arithmetic. Whether it survives quantisation is precisely the kind of
question Phase~1 answers (\cref{ch:codesign}). One edge case: $q=0$ means
$\operatorname{logit}(q)=-\infty$ and $p=1$, which has to be handled explicitly.
\end{conceptbox}

% =========================================================================
\section{Memory: the dominant cost}\label{sec:workload:memory}

The recursive estimators carry state from one frame to the next. Introspecting
a running \code{SpeechEnhancer} gives the inventory in
\cref{tab:workload:state} (measured).

\begin{table}[htbp]
\centering\small
\caption{Arrays that persist between frames (measured by introspection).}
\label{tab:workload:state}
\begin{tabularx}{\textwidth}{@{}llrL@{}}
\toprule
Array & Shape & words/bin & Role \\
\midrule
\code{\_tracker.\_store}         & $8\times1025$ & 8 & subwindow minima, tracker 1 \\
\code{\_tracker.\_running}       & $1025$ & 1 & running minimum, tracker 1 \\
\code{\_tracker\_tilde.\_store}  & $8\times1025$ & 8 & subwindow minima, tracker 2 \\
\code{\_tracker\_tilde.\_running}& $1025$ & 1 & running minimum, tracker 2 \\
\code{\_s}, \code{\_s\_tilde}    & $1025$ each & 2 & smoothed periodograms $S$, $\tilde S$ \\
\code{\_lambda\_d}               & $1025$ & 1 & noise PSD estimate \\
\code{\_prev\_clean\_power}, \code{\_xi} & $1025$ each & 2 & Wiener decision-directed memory \\
\code{\_norm\_freq}              & $1025$ & (1) & constant ramp for tilt --- ROM or computed, not state \\
\midrule
\multicolumn{2}{@{}l}{\textbf{Mutable state}} & \textbf{23} & $23\times1025=\num{23575}$ words \\
\midrule
\code{window}, \code{\_in\_buffer}, \code{\_out\_buffer} & $3\times1024$ & & window ROM, input and overlap-add buffers \\
\bottomrule
\end{tabularx}
\end{table}

\begin{figure}[htbp]
\centering
\begin{tikzpicture}
\begin{axis}[
  xbar stacked, width=0.92\textwidth, height=34mm,
  bar width=7mm, xmin=0, xmax=24,
  ytick=\empty, xlabel={words per bin (measured)},
  xlabel style={font=\sffamily\footnotesize}, tick label style={font=\sffamily\footnotesize},
  legend style={at={(0.5,1.05)},anchor=south,legend columns=-1,font=\sffamily\scriptsize,draw=none},
  axis line style={draw=conceptc}, enlarge y limits=0.6,
]
\addplot[fill=pitc!70,draw=white] coordinates {(16,0)};
\addplot[fill=pitc!35,draw=white] coordinates {(2,0)};
\addplot[fill=brand!60,draw=white] coordinates {(3,0)};
\addplot[fill=whyc!60,draw=white] coordinates {(2,0)};
\legend{tracker subwindow minima (16),tracker running minima (2),{IMCRA $S,\tilde S,\lambda_d$ (3)},Wiener memory (2)}
\end{axis}
\end{tikzpicture}
\caption{Where the 23 words of per-bin state go: 18 of them are inside the two
minimum trackers.}
\label{fig:workload:state}
\end{figure}

At \num{23575} words this is \qty{46.0}{\kibi\byte} at 16~bits,
\qty{69.1}{\kibi\byte} at 24~bits and \qty{92.1}{\kibi\byte} at 32~bits (measured).
That is a lot of on-chip memory for a hearing aid. Reading and writing SRAM also
typically costs more energy than the arithmetic done on the data, so memory is
likely to dominate both area and energy. Three observations shape Phase~1:

\begin{enumerate}
  \item \textbf{The trackers are the target.} 18 of the 23 words per bin hold
  minima. A minimum is unchanged by any monotone transform, so the trackers can
  store coarse logarithmic values (O5). As an illustration (not yet a result):
  storing the 18 tracker words at 8~bits and the other 5 at 16~bits gives
  $18\times8+5\times16=224$ bits per bin, against $23\times24=552$. That is a
  2.5$\times$ reduction before you change anything else.
  \item \textbf{The number of bins multiplies everything.} Removing the 2$\times$
  zero padding (O6) halves the bins from 1025 to 513. Moving to \qty{16}{\kHz}
  with a 128-sample window leaves 65 bins (\cref{sec:workload:latency}). With
  the 224-bit layout above, 65 bins need about \qty{1.8}{\kibi\byte} in total.
  \item \textbf{Memory \emph{accesses} matter as much as memory size.} As
  written, each tracker reads all 9 stored values for every bin in every frame
  (\file{imcra.py:51}), even though the 8 subwindow minima only change once every
  $V=15$ frames. If you cache their minimum in one extra word, the per-frame
  reads drop to 2 per tracker. The full 8-word scan is then needed only when a
  slot is overwritten. That is about $2+8/15\approx2.5$ reads per frame instead
  of 9, paid for with 1 extra word per bin per tracker. This kind of trade
  (slightly more storage for far fewer accesses) appears constantly in low-power
  design.
\end{enumerate}

\begin{tipbox}[Pre-roll needs no buffer]
The thesis stores 5 whole pre-roll frames and then takes their per-bin minimum
(\file{imcra.py}, \code{seed}). Since only the minimum is used, hardware can keep
a running minimum in the $\lambda_d$ state memory that is about to be seeded
anyway. That costs 0 extra words instead of $5\times1025$.
\end{tipbox}

% =========================================================================
\section{Latency anatomy}\label{sec:workload:latency}

\Cref{fig:workload:latency} shows why every sample is delayed by exactly one
window length. A block of 512 samples must be fully captured before its frame can
be analysed. The overlap-add output for that block is only complete once the
\emph{next} frame has been added in. The total is 1024 samples, or
\qty{21.33}{\ms} at \qty{48}{\kHz} (measured), before any converter or operating
system buffering. Hearing-aid research generally treats total delays above
roughly \qty{10}{\ms} as problematic, and open fittings are less tolerant still,
because the direct and processed sound interfere and cause comb
filtering~\cite{stone1999delay,stone2008delay}.

\begin{figure}[htbp]
\centering
\begin{tikzpicture}[font=\sffamily\scriptsize,x=6mm,y=7mm]
  \draw[-{Stealth},conceptc] (0,0) -- (21,0) node[right]{time};
  \foreach \t/\l in {0/0,5/10.67,10/21.33,15/32}{
    \draw[conceptc] (\t,0.1) -- (\t,-0.1) node[below]{\l\,ms};}
  % input blocks
  \node[anchor=east] at (-0.2,2.2) {input};
  \draw[fill=brand!25,draw=brand] (0,1.9) rectangle (5,2.5) node[midway]{block $k{-}1$};
  \draw[fill=brand!10,draw=brand] (5,1.9) rectangle (10,2.5) node[midway]{block $k$};
  \draw[fill=brand!5,draw=brand,dashed] (10,1.9) rectangle (15,2.5) node[midway]{block $k{+}1$};
  % processing
  \node[anchor=east] at (-0.2,1.15) {process};
  \draw[fill=pitc!20,draw=pitc] (10,0.9) rectangle (11.3,1.4);
  \node[anchor=west,text=pitc] at (11.4,1.15) {frame $[k{-}1,k]$: FFT, gains, IFFT, OLA};
  % output
  \node[anchor=east] at (-0.2,3.2) {output};
  \draw[fill=donec!25,draw=donec] (10,2.9) rectangle (15,3.5) node[midway]{block $k{-}1$ (enhanced)};
  \draw[{Stealth}-{Stealth},pitc,line width=0.7pt] (0,3.9) -- node[above]{algorithmic latency = 1024 samples = 21.33 ms} (10,3.9);
\end{tikzpicture}
\caption{Why the thesis pipeline delays every sample by one full window. Block
$k-1$ cannot leave until block $k$ has arrived, because its overlap-add output
needs the next frame. Processing time, drawn short here, adds to this.}
\label{fig:workload:latency}
\end{figure}

For a frame-based design with 50\,\% overlap, algorithmic latency equals the
window length. Latency is therefore a design variable you can trade against
frequency resolution and memory (\cref{tab:workload:configs}, measured state
sizes). This is research question RQ2.

\begin{table}[htbp]
\centering\small
\caption{Alternative frame configurations (no zero padding, hop = window/2).}
\label{tab:workload:configs}
\begin{tabular}{@{}rrrrrr@{}}
\toprule
$f_s$ & window & bins & latency & frame rate & state (words) \\
\midrule
\qty{48}{\kHz} & 2048 FFT$^\ast$ & 1025 & \qty{21.33}{\ms} & \qty{93.75}{\Hz} & \num{23575} \\
\qty{48}{\kHz} & 1024 & 513 & \qty{21.33}{\ms} & \qty{93.75}{\Hz} & \num{11799} \\
\qty{48}{\kHz} & 512  & 257 & \qty{10.67}{\ms} & \qty{187.5}{\Hz} & \num{5911} \\
\qty{16}{\kHz} & 256  & 129 & \qty{16.00}{\ms} & \qty{125}{\Hz} & \num{2967} \\
\qty{16}{\kHz} & 128  & 65  & \qty{8.00}{\ms}  & \qty{250}{\Hz} & \num{1495} \\
\qty{16}{\kHz} & 64   & 33  & \qty{4.00}{\ms}  & \qty{500}{\Hz} & \num{759} \\
\bottomrule
\multicolumn{6}{@{}l}{\footnotesize $^\ast$ The thesis configuration: 1024-sample window zero-padded to 2048.}
\end{tabular}
\end{table}

\begin{pitfallbox}[Time constants are defined in frames]
IMCRA's smoothing factors ($\alpha_s=0.9$, $\alpha_d=0.85$), the
decision-directed factor (0.98) and the minimum-search window ($D=U\times V=120$
frames, about \qty{1.28}{\s}) are all expressed \emph{per frame}. If you halve the
hop and keep these values, every estimator suddenly adapts twice as fast. A
frame-configuration study that forgets this compares different algorithms, not
different frame sizes. To keep the time constants fixed when the hop changes
from $H$ to $H'$, rescale them: $\alpha' = \alpha^{H'/H}$, and $D' = D\,H/H'$
frames. \Cref{ch:codesign} builds this into the sweep.
\end{pitfallbox}

% =========================================================================
\section{The control plane}\label{sec:workload:control}

The intent path runs rarely, but each run is large (\cref{tab:workload:control}).
Almost all of the cost is in the sentence encoder. The classifier and severity
heads are tiny matrix--vector products. This is a direct consequence of your
thesis finding that the classes are linearly separable in the embedding space.

\begin{table}[htbp]
\centering\footnotesize
\caption{Control-plane cost per utterance. Encoder figures are estimates
($2\times$ non-embedding parameters $\times$ about 14 tokens; the complaint
sentences average 7.8 words). Head figures are exact from the model shapes.}
\label{tab:workload:control}
\begin{tabularx}{\textwidth}{@{}lrlL@{}}
\toprule
Component & Parameters & Compute per utterance & Weights (fp32 / int8) \\
\midrule
\code{all-mpnet-base-v2}~\cite{song2020mpnet,reimers2019sbert} & $\approx$110\,M & $\approx$2.4\,GFLOP & $\approx$440\,MB / $\approx$110\,MB \\
\code{all-MiniLM-L6-v2}~\cite{wang2020minilm} & $\approx$22.7\,M & $\approx$0.3\,GFLOP & $\approx$90\,MB / $\approx$23\,MB \\
\midrule
Linear SVC head (one-vs-one) & $15\times768$ & \num{11520} MAC + voting & \qty{45}{\kibi\byte} / \qty{11.25}{\kibi\byte} \\
Logistic regression head & $6\times768$ & \num{4608} MAC + argmax & \qty{18}{\kibi\byte} / \qty{4.5}{\kibi\byte} \\
Ridge severity head & $768$ & 768 MAC + clip & \qty{3}{\kibi\byte} / \qty{0.75}{\kibi\byte} \\
\bottomrule
\end{tabularx}
\end{table}

\begin{whybox}[A detail that matters: the selected classifier is one-vs-one]
The thesis selected \code{SVC(kernel="linear", probability=True)}
(\file{src/intent_se/nlp/classifier.py}). scikit-learn's \code{SVC} trains one
binary classifier for \emph{every pair} of classes: 15 hyperplanes for 6 classes,
so 15 rows rather than 6. Its probability output, which the confidence gate
uses, comes from Platt scaling followed by an iterative pairwise-coupling
solve~\cite{platt1999probabilistic,wu2004pairwise}. That is awkward to put in
hardware. Logistic regression tied the SVM on cross-validated macro-F1 (0.964).
It needs 2.5$\times$ fewer multiply--accumulates, and its confidence is simply
its largest score. The thesis broke the tie on generalisation, which was right
for a laptop. On a chip, the tie may break the other way. Phase~1 measures this
(opportunity O9).
\end{whybox}

% =========================================================================
\section{Decisions that flip under hardware constraints}\label{sec:workload:flips}

Several thesis decisions were correct for a laptop proof of concept but should
be revisited for hardware (\cref{tab:workload:flips}). Explaining these flips
clearly is one of the most valuable things you can say in an interview, because
it shows you understand that \enquote{optimal} depends on the cost model.

\begin{table}[htbp]
\centering\footnotesize
\caption{Thesis decisions to revisit, and where in this guide each is tested.}
\label{tab:workload:flips}
\begin{tabularx}{\textwidth}{@{}LLLl@{}}
\toprule
Thesis decision & Thesis reason & Under hardware constraints & Tested in \\
\midrule
mpnet (110\,M) over MiniLM & \enquote{Inference speed is not a bottleneck} & The encoder is the dominant control-plane cost & RQ4, \cref{ch:mlaccel} \\
2$\times$ zero padding (2048-pt FFT) & Smoother gain curves & Doubles FFT work and per-bin state & H2, \cref{ch:codesign} \\
\qty{48}{\kHz}, 1024 window & Soundcard block size, frequency resolution & \qty{21.33}{\ms} latency, 1025 bins of state & RQ2, \cref{ch:codesign} \\
float64 everywhere & Free in Python & Word length sets area, energy and memory & RQ1, \cref{ch:codesign} \\
Linear SVM (one-vs-one, Platt) & Best val$\to$test generalisation & LR ties on F1, needs 2.5$\times$ fewer MACs and has a simple confidence & H3, \cref{ch:codesign,ch:mlaccel} \\
Tilt computed per bin per frame & Simple code & Precompute at control rate (O2) & \cref{ch:codesign,ch:rtl} \\
Store 5 pre-roll frames & Simple code & Running minimum in place, no buffer & \cref{ch:rtl} \\
\bottomrule
\end{tabularx}
\end{table}

% =========================================================================
\section{Summary of hardware opportunities}\label{sec:workload:opportunities}

The rest of the guide refers to these opportunities by label
(\cref{tab:workload:opps}). Each one is a hypothesis about cost and accuracy,
not a guaranteed win. Phase~1 measures the accuracy side and Phases~2--5
measure the cost side.

\begin{table}[htbp]
\centering\footnotesize
\caption{Hardware opportunities found in the thesis code.}
\label{tab:workload:opps}
\begin{tabularx}{\textwidth}{@{}lLLl@{}}
\toprule
ID & Idea & What it saves & Validated in \\
\midrule
O1 & Gain LUT $G(\xi_\text{dB})=\max\big((\xi/(1+\xi))^\beta,G_\text{floor}\big)$, rebuilt by the control plane, double-buffered & a \code{pow} and a division per bin; turns the NLP$\to$DSP interface into a table write & \cref{ch:codesign,ch:rtl} \\
O2 & Tilt as a geometric accumulator $g_k=g_\text{out}\,r^k$, $r=10^{\text{tilt}/(20(K-1))}$ & a $10^x$ per bin & \cref{ch:codesign,ch:rtl} \\
O3 & Cache $H=\xi/(1+\xi)$ from the Wiener stage for next frame's IMCRA & a reciprocal per bin & \cref{ch:codesign} \\
O4 & Share reciprocals across the 9 distinct denominators; LUT seed + Newton--Raphson & 3 of 12 divisions, then cheap reciprocals & \cref{ch:codesign,ch:rtl} \\
O5 & Log-domain values in the trackers and SNR comparisons & bits in 18 of 23 state words; divisions become subtractions & \cref{ch:codesign} \\
O6 & Remove the 2$\times$ zero padding & half the FFT work and half the state & \cref{ch:codesign} \\
O7 & $e^{-\nu}$ by LUT or piecewise-linear (or the sigmoid rewrite, \cref{eq:workload:sigmoid}) & an \code{exp} per bin & \cref{ch:codesign} \\
O8 & Block-floating-point FFT and power computation & word length in the FFT & \cref{ch:codesign,ch:rtl} \\
O9 & Logistic-regression or margin-based confidence; int8 head with Ridge folded in as a 7th row & 2.5$\times$ fewer MACs; no iterative probability solve & \cref{ch:codesign,ch:mlaccel} \\
\bottomrule
\end{tabularx}
\end{table}

% =========================================================================
\section{Do it yourself}\label{sec:workload:diy}

\begin{stepbox}[label=st:workload:reproduce]{Reproduce this chapter's measurements}
From the root of the thesis repository run:
\begin{lstlisting}[style=shell]
PYTHONPATH=src python docs/research_guide/scripts/measure_workload.py
\end{lstlisting}
\textbf{Check:} the output reports 23 mutable words per bin (18 in the trackers),
\num{23575} words of state, a frame period of \qty{10.667}{\ms} and a latency of
\qty{21.333}{\ms}. These numbers must match \cref{app:factsheet} exactly. The
Python time per frame varies from machine to machine.

\textbf{Effort:} 15 minutes.
\end{stepbox}

\begin{stepbox}[label=st:workload:opcount]{Turn the operation estimate into a measurement (optional)}
\Cref{tab:workload:ops} was counted by hand. You can count automatically
instead. Wrap the numpy arrays inside \code{IMCRA.update} and
\code{ParametricWienerFilter.compute\_gain} in a thin \code{ndarray} subclass
whose \code{\_\_array\_ufunc\_\_} increments a counter per ufunc
(\code{multiply}, \code{divide}, \code{minimum}, \code{exp}, \ldots) by the
number of elements, then process one frame. Divide the counts by 1025 to get
per-bin numbers.

\textbf{Check:} your divisions per bin equal 12 and your \code{exp}/\code{power}
calls equal 2. If not, find the discrepancy. It will teach you something about
the code.

\textbf{Effort:} 2--3 hours. Put the result in your lab notebook. It is a good
first entry.
\end{stepbox}

\begin{donebox}[You are ready for the project chapters when]
\begin{checklist}
  \item You can explain, without notes, why memory rather than arithmetic speed
        is the main hardware problem for this system.
  \item You can name the two most expensive nonlinear operations per bin and how
        each will be replaced.
  \item You know where the \qty{21.33}{\ms} of latency comes from and what
        changing the window does to latency, memory and IMCRA time constants.
  \item You have reproduced the measured numbers on your own machine.
\end{checklist}
\end{donebox}
